% MUST use a4paper option
% MAY use twoside, smaller font, and other class - but not for Självständigt arbete i IT
\documentclass[a4paper,12pt]{article}
% Use UTF-8 encoding in input files
\usepackage[utf8]{inputenc}
% Use T1 font encoding to make the \hyphenation command work with UTF-8
\usepackage[T1]{fontenc}

% Om ni skriver på svenska, använd denna rad:
\usepackage[english,swedish]{babel}
% If you are writing in English, use the following line INSTEAD of the previous (note order of parameters):
% \usepackage[swedish,english]{babel}

% Use the template for thesis reports
\usepackage{UppsalaExjobb}

% För att göra ett index behövs
%  - \usepackage{makeidx}
%  - \makeindex i "preamble", dvs före \begin{document}
%  - \printindex, typiskt sist, före \end{document}
% - och att man lägger in \index{ord} på olika ställen i dokumentet
\usepackage{makeidx}
\makeindex


% Designval: per default används styckesindrag, men ibland blir det snyggare/mer lättläst med tomrad mellan stycken. Det åstadkoms av de följande raderna.
% Tycker ni om styckesindrag mera, kommentera bort nästa två rader.
\parskip=0.8em
\parindent=0mm

% Designval: vill ni ha en box runt figurer istället för strecken som är default, av-kommentera raden nedan. Obs att både \floatstyle och \restylefloat behövs.
%\floatstyle{boxed} \restylefloat{figure}

\begin{document}
% För att ställa in parametrar till IEEEtranS/IEEEtranSA behöver detta ligga här (före första \cite).
% Se se IEEEtran/bibtex/IEEEtran_bst_HOWTO.pdf, avsnitt VII, eller sista biten av IEEEtran/bibtex/IEEEexample.bib.
%%%% OBS: här ställer ni t.ex. in hur URLer ska beskrivas.
\bstctlcite{rapport:BSTcontrol}

% Set title, and subtitle if you have one
\title{Rapportmall för självständigt arbete} % och uppsatsmetodik
% Use this if you have a subtitle
%\subtitle{beskrivande men gärna lockande}
\subtitle{version vinter 2022}

% Set author names, separated by "\\ " (don't forget the space, or use newline)
% List authors alphabetically by LAST NAME (unless someone did significantly more/less, which should not be the case)
% For drafts, include your email addresses to make it easier to send peer reviews
\author{Sofia Cassel \\ Björn Victor (bjorn.victor@it.uu.se)}

% Visa datum på svenska på förstasidan, även om ni skriver på engelska!
\date{\begin{otherlanguage}{swedish}  %\foreignlanguage doesn't seem to affect \today?
\today
\end{otherlanguage}}

% Använd detta om året för rapporten inte är innevarande år
%\year=2018

% Ange handledare, ämnesgranskare, examinator om dessa finns
% Extern handledare: t.ex på företag ni arbetat med?
\exthandledare{Hand Ledare, Firma Ment AB}
% Intern(a) handledare, i bokstavsordning på efternamnet
\handledare{NN och Björn Victor}

% This creates the title page
\maketitle

% Change to frontmatter style (e.g. roman page numbers)
\frontmatter

%%%% OBS: Läs också källkoden till alla instr-X.tex.
%%%% Tips: ni kan använda separata filer för de olika delarna i er rapport på motsvarande sätt,
%%%% men använd inte samma filnamn!

\begin{abstract}
\input{instr-abstract}
\end{abstract}

\begin{sammanfattning}
\input{instr-sammanfattning}
\end{sammanfattning}

% Innehållsförteckningen här.
\tableofcontents

% Här kan man också ha \listoffigures, \listoftables

\cleardoublepage

%%%%%%%%%%%%%%%% Ta bort allt mellan här och \mainmatter (inkl \newpage) (men inte \mainmatter) i slutversionen
\input{instr-howto-use}
\newpage
%%%%%%%%%%%%%%%% OBS! Ta bort allt mellan \mainmatter och här (inkl \newpage) i slutversionen

% Change to main matter style (arabic page numbers, reset page numbers)
\mainmatter

% Here comes the text of the report.

\section{Introduktion eller Inledning / Introduction}
\label{sec:introduktion}
With the rapid digitalization of real-world systems, new challenges arise. One such example is the transition to digital cashier systems. While this shift offers retailers numerous advantages, such as increased effiency, accuracy and scalability, it also introduceses risks that must be properly managed. This project, named AI that catches mistakes, aims to explores how user errors and anomalies in cashier transactions can be detected through the use of AI. The system makes use of AI to flag mistakes made by cashiers in real time, allowing them to be corrected.

\paragraph{Tillkännagivande eller Tack / Acknowledgement}

\paragraph{Redovisning av arbetsfördelning / Declaration of division of labor}


\section{Bakgrund / Background}
\label{sec:bakgrund}
Shrinkage is a term used to describe a loss of inventory, and in the case of retail it usually occurs in the Point Of Sale (POS), in other words the time and place the transaction occurs. Shrinkage can be caused by several factors, including theft, administrative errors, fraud, or mistakes made during checkout. Any kind of manipulation of the cashier system, whether intentional or unintentional, is referred to as a POS exception. This term encompasses forms of internal shrinkage, such as distributions of unauthorized coupons, falsifying prices, and other internal actions that lead to shrinkage.\footnote{\cite{netsuite-shrink}} According to an article by Tom Howkings,a retail transformation architect at Deloitte, 40\% of retail theft is done by employees internally within the company, and is believed to be under reported.\footnote{\cite{deloitte-shrink}} Additionally, the National Retail Federation estimated that shrinkage led American retailers to lose \$112 billion in 2022, creating an average shrink rate of 1.6\%, meaning the average percent of inventory that disappears due to shrinkage.

The traditional solution to minimize shrinkage has been to strengthen existing routines regarding inventory audits, employee training, and manual oversight. Retailers often rely on stock counts, surveillance systems, and revising transactional data to identify discrepancies. While these methods can be effective, they are typically reactive rather than proactive, meaning that losses are often discovered only after they have already occurred. Additionally, manual processes are time-consuming, prone to human error, and difficult to scale in large retail environments with high transaction volumes.

\section{Syfte, mål, och motivation / Purpose, aims, and motivation}\label{sec:syfte}
As mentioned in the introduction, the goal of this work is to develop a system that can effectively flag and catch mistakes in real-time in cashier transactions. The purpose of this is to reduce shrinkage caused by cashiers, which accounts for major economic losses globally in retail. Introducing a system that relies on AI to address this could lead to decreased resources spent on manual processes that are prone to human error. Additionally, this would create a proactive approach to decreasing shrinkage, rather than a reactive one.

\subsection{Avgränsningar / Delimitations}
As for the limitations, the system is limited by the amount and type of data available. There could occur different kinds of mistakes and manipulation in the system at different types of stores, leading the system to perform worse if used in cases it is not trained for. It should also be noted that the term mistakes is broad, and encompasses mistakes that don't show up in the transaction log, such as not scanning an item. Therefore the system can't identify or flag these types of mistakes, since it has no way to detect it through the transaction log.

\section{Relaterat arbete / Related work}

\section{Teori / Theory}
\label{sec:teori}

\section{Metod/Tillvägagångssätt/Tekniker eller Method/Approach/Techniques}
\label{sec:metod}

\section{Systemstruktur / System structure}
\label{sec:systemstruktur}

\section{Krav och utvärderingsmetoder / Requirements and evaluation methods}
\label{sec:krav}

\section{DEL x: Implementation av XYZ}
\label{sec:del-x}

\section{Utvärderingsresultat / Evaluation results}

\section{Resultat och diskussion / Results and discussion}
\label{sec:resultat}

\section{Framtida arbete / Future work}
\label{sec:framtida}

\section{Slutsatser / Conclusions}
\label{sec:slutsatser}

%%%% Referenser - SE OCKÅ APPENDIX

% Use one of these:
%   IEEEtranS gives numbered references like [42] sorted by author,
%   IEEEtranSA gives ``alpha''-style references like [Lam81] (also sorted by author)
%\bibliographystyle{IEEEtranS}
\bibliographystyle{IEEEtranSA}

% Here comes the bibliography/references.
% För att göra inställningar för IEEEtranS/SA kan man använda ett speciellt bibtex-entry @IEEEtranBSTCTL,
% se IEEEtran/bibtex/IEEEtran_bst_HOWTO.pdf, avsnitt VII, eller sista biten av IEEEtran/bibtex/IEEEexample.bib.
\bibliography{bibconfig,refs}
%\bibliography{refs}

\newpage
\appendix %%%% markerar att resten är appendixar
%%%% I er egen version, ta bort allt nedan (utom \end{document})
\input{instr-appendix}

% Om ni har ett index
\makeatletter
\renewenvironment{theindex}
               {\if@twocolumn
                  \@restonecolfalse
                \else
                  \@restonecoltrue
                \fi
                \twocolumn[\section{\indexname}]%
                \@mkboth{\MakeUppercase\indexname}%
                        {\MakeUppercase\indexname}%
                \thispagestyle{plain}\parindent\z@
                \parskip\z@ \@plus .3\p@\relax
                \columnseprule \z@
                \columnsep 35\p@
                \let\item\@idxitem}
               {\if@restonecol\onecolumn\else\clearpage\fi}
\makeatother
\printindex
\end{document}
